%% ECC specific appendix

%%% ECC curve specifikus
\section{Basics for Elliptic Curve Theory}
\label{sec:apx.ecc.curves}

% 2 curve csalad:
An elliptic curve is the set of solutions to the \textit{Weierstrass form}
\textit{equation \ref{eq:ecc.weierstrass}} with respect to
$(x,y) \in \mathcal{F}_{q^m} \times \mathcal{F}_{q^m}$.\cite{ecc.fam.param}

% Weierstrass equation
\begin{equation}
\label{eq:ecc.weierstrass}
E: y^2 = x^3 + ax + b,
\end{equation}
where $a,b \in \mathcal{F}_{q^m}$ $\Delta \neq 0 $ (discriminant)
and a point at infinity ($\infty$).\cite{ecc.fam.param}

Also $E'$ a degree $d$ twist of $E$ is an isomorphic curve to $E$
over the algebraic closure of $\mathcal{F}_{q^m}$. The only possible
degrees for elliptic curves are $d \in {2,3,4,6}$.\cite{ecc.fam.param}

%   TODO: BF01 alapjan megnezni, majd, hogy kell e itt ez az elmeleti resz.
%   TODO: A pairinghez legyen library hasznalva a side channel, timing, etc.
%           attack miatt --> de az elemeleti reszben be kell irni, hogy mi miert
%           lett hasznalva
\subsubsection{Hasse condition}

The order $n$ of the curve is the number of points that satisfy
the curve equation. The \textit{Hasse conddition} states the following
\textit{\ref{eq:ecc.hassecond} equation}:

\begin{equation}
    \label{eq:ecc.hassecond}
    n = q^m +1 - t,~~ \lvert t \rvert \leq 2 \sqrt{q^m}
\end{equation}

We call the curve \textit{supersingular} if $q$ divides $t$.

The order of point $P$ is the smallest integer $r$ such that
$rP = \infty$. It is also true that the order of the point
is a dividant of the curve order: $r|n$.\cite{ecc.fam.param}

The $r$-torsion subgroup (denoted as $E(\mathcal{F}_{q^m}[r])$) is the
set of points $P$ in which their oder divides $r$.\cite{ecc.fam.param}

%   TODO: ETSI alapjan ha mas curve-re is ellenorznik
% Edwards Curve
% Montgomery curve

\subsection{Bilinear pairings}
\label{ssec:apx.ecc.curves.pairings}

Let $\mathcal{G}_1 = \left< P \right>$ and $\mathcal{G}_2 = \left< P \right>$
be additive groups and $\mathcal{G}_T$ be a
multiplicative group such that $\left\| \mathcal{G}_1 \right\| = \left\| \mathcal{G}_2 \right\| = \left\| \mathcal{G}_T \right\|= \mathrm{prime}~ r$

We call the $e: \mathcal{G}_1 \times \mathcal{G}_2 \rightarrow \mathcal{G}_T $
addmissible biilinear map if \textit{bilinearity} and
\textit{non-degeneracy} properties are satisfied:

\begin{itemize}
    \item[] \textbf{Bilinearity: } given $(V, W) \in \mathcal{G}_1 \times \mathcal{G}_2$ and
    $(a,b) \in \mathcal{Z}^{*}_{r}$:
    \begin{itemize}
        % Azonossagok:
        \item[] $e(aV, bW) = e(V, W)^{ab} = e(abV, W) = e(V, abW) = e(bV, aW)$
        \end{itemize}
    \item[] \textbf{Non-degeneracy:} $e(P,Q) \neq 1_{\mathcal{G}_T}$, 
    \end{itemize}

% Properties of general pairings
%% TODO: Ide a konkret weil pairinget hozzarakni +
%%       a unit teszteket a valasztott parameterek alapjan ide berakni
\subsubsection{Properties of pairings}

In case of pairings ( $e: \mathcal{G}_1 \times \mathcal{G}_2 \rightarrow \mathcal{G}_T $) there are some properties we can rely on.
Namely, $\mathcal{G}_1$ is typically a subgroup of $E(\mathcal{F}_q)$,
$\mathcal{G}_2$ is typically a subgroup of $E(\mathcal{F}_q^k)$,
lastly $\mathcal{G}_T$ is a multiplicative subgroup of $\mathcal{F}^*_{q^k}$.\cite{ecc.fam.param}

% In practice we want small `k` for computable pairing
The embedding degree of $\mathcal{F}_q$ is the smallest integer $k$,
in which we can embed the $r$-torsion group: $r|(q^k-1)$. For efficiency reasons
we want to the largest $d$ such that $d|k$.\cite{ecc.fam.param}
The families of ordinary curves usually have $k < 50$, random curves have
$k \approx q$, but for supersingular curves we have embedding degree of
$k \leq 6$.\cite{ecc.fam.param}

From the \textit{Adversary}'s point of view $k$ is usually very large and hence mapping
the problem into the $\mathcal{F}_{p^k}$ field is useless for an attacker.
As for super-singular curves that the $k$ embedding degree
is small, but on the other hand, If we choose the curves large enough,
discrete log will still be hard despite the attacks.\cite{crypto.ivan.cryptology}




For this practical reasons we choosed standardized curves mentioned in
\textit{subsection \label{ssec:ibe.bonehfranklin.scheme}}.

